% ============================================================
%  Hyperlocal Air Quality Mapping in Chicago
%  A Citizen Science Project in Urban Aerosol Transport
%  Project Overview Document
% ============================================================
%  Compile with:  pdflatex chicago_run_bike_monitor.EDITED.TEX
%  (run twice to resolve cross-references)
%  EDIT [mechanical, 2026-07-22]: AUDIT_papers.md Paper-3 issue #5 -- this line named a
%  file "chicago_airquality_project.tex" that does not exist in the tree. Corrected.
% ============================================================

\documentclass[11pt, letterpaper]{article}

% ── packages ────────────────────────────────────────────────
\usepackage[margin=1in]{geometry}
\usepackage[T1]{fontenc}
\usepackage[utf8]{inputenc}
\usepackage{lmodern}
\usepackage{microtype}
\usepackage{parskip}          % space between paragraphs, no indent
\usepackage{titlesec}         % custom heading styles
\usepackage[table]{xcolor}
\usepackage{booktabs}
\usepackage{array}
\usepackage{enumitem}
\usepackage{tcolorbox}
\usepackage{fancyhdr}
\usepackage{hyperref}
\usepackage{amsmath}
\usepackage{amssymb}
\usepackage{ragged2e}
\tcbuselibrary{skins, breakable}

% ── colours ─────────────────────────────────────────────────
\definecolor{navy}{HTML}{1B3A6B}
\definecolor{teal}{HTML}{1D6A72}
\definecolor{gold}{HTML}{C8860A}
\definecolor{lightgray}{HTML}{F2F4F7}
\definecolor{midgray}{HTML}{D0D4DB}
\definecolor{lteal}{HTML}{EAF4F5}
\definecolor{darktext}{HTML}{1A1A1A}

% ── heading styles ───────────────────────────────────────────
\titleformat{\section}
  {\color{navy}\Large\bfseries}
  {\thesection.}{0.6em}{}
  [\vspace{2pt}\color{teal}{\titlerule[0.8pt]}\vspace{4pt}]

\titleformat{\subsection}
  {\color{teal}\large\bfseries}
  {\thesubsection}{0.6em}{}

\titleformat{\subsubsection}
  {\color{gold}\normalsize\bfseries}
  {\thesubsubsection}{0.6em}{}

\titlespacing{\section}{0pt}{18pt}{6pt}
\titlespacing{\subsection}{0pt}{12pt}{4pt}
\titlespacing{\subsubsection}{0pt}{8pt}{3pt}

% ── hyperref ────────────────────────────────────────────────
\hypersetup{
  colorlinks = true,
  linkcolor  = navy,
  urlcolor   = teal,
  citecolor  = teal,
  pdftitle   = {Hyperlocal Air Quality Mapping in Chicago},
  pdfauthor  = {Northwestern University PhD Project},
}

% ── callout box ─────────────────────────────────────────────
\tcbset{
  callout/.style={
    enhanced,
    breakable,
    colback=lteal,
    colframe=teal,
    fonttitle=\bfseries\color{navy},
    boxrule=0.7pt,
    arc=3pt,
    left=8pt, right=8pt, top=6pt, bottom=6pt,
    before upper={\setlength{\parskip}{3pt}},
  }
}

% ── table helpers ────────────────────────────────────────────
\newcolumntype{L}[1]{>{\RaggedRight\arraybackslash}p{#1}}
\newcolumntype{C}[1]{>{\Centering\arraybackslash}p{#1}}

\renewcommand{\arraystretch}{1.3}

% ── header / footer ─────────────────────────────────────────
\pagestyle{fancy}
\fancyhf{}
\renewcommand{\headrulewidth}{0.4pt}
\renewcommand{\headrule}{\hbox to\headwidth{\color{midgray}\leaders\hrule height \headrulewidth\hfill}}
\fancyhead[L]{\small\color{teal} Hyperlocal Air Quality Mapping in Chicago}
\fancyhead[R]{\small\color{midgray} Page \thepage}

% ── list style ───────────────────────────────────────────────
\setlist[itemize]{
  leftmargin=1.4em,
  itemsep=2pt,
  parsep=0pt,
  topsep=4pt,
}
\setlist[enumerate]{
  leftmargin=1.6em,
  itemsep=2pt,
  parsep=0pt,
  topsep=4pt,
}

% ============================================================
\begin{document}
% ============================================================

% ── cover page ───────────────────────────────────────────────
\thispagestyle{empty}
\begin{center}
  \vspace*{3cm}
  {\color{navy}\Huge\bfseries Hyperlocal Air Quality\\[6pt] Mapping in Chicago}\\[16pt]
  {\color{teal}\Large\itshape A Citizen Science Project in Urban Aerosol Transport}\\[20pt]
  {\color{gold}\rule{0.5\textwidth}{1.2pt}}\\[20pt]
  {\large PhD Research Project \quad|\quad Northwestern University}\\[6pt]
  {\normalsize Department of Physics \quad|\quad Soft Condensed Matter \& Fluid Dynamics}\\[6pt]
  {\normalsize 2026\,--\,2027}
\end{center}
\vfill
\begin{center}
  \small\color{midgray}
  Hardware instrumentation $\cdot$ Open dataset $\cdot$ Urban aerosol physics $\cdot$ Citizen science
\end{center}
\newpage

% ── table of contents ────────────────────────────────────────
\tableofcontents
\newpage

% ============================================================
\section{Project Summary}
% ============================================================

Chicago has just five EPA air quality monitoring stations for a city of 2.7 million people---a spatial
resolution of roughly one sensor per 540{,}000 residents. This project deploys a lightweight,
wearable sensor package on everyday running and cycling routes to collect fine-grained measurements
of particulate matter (PM\textsubscript{2.5}, PM\textsubscript{10}), volatile organic compounds
(VOCs), temperature, humidity, and pressure, each paired with a GPS coordinate and timestamp.

The result will be an open dataset of street-level air quality at spatial resolutions of 5--10\,m
and time resolutions of 1\,s---roughly 100$\times$ finer than anything currently publicly available
for Chicago. On top of that dataset, the project applies established fluid mechanics and atmospheric
physics models (Gaussian plume dispersion, street canyon recirculation, Stokes-number particle
sorting, lake breeze dynamics) to turn casual outdoor exercise into a genuine scientific
investigation of urban air quality.

This is a solo, end-to-end project: you design the hardware, collect every data point yourself,
build the analysis pipeline, and tell the final scientific story---from sensor to published result.

\newpage

% ============================================================
\section{Scientific Goals and Hypotheses}
% ============================================================

Each hypothesis below is specific, directional, and testable with the dataset you will collect.
They are ordered from most tractable (large expected signal, few confounders) to most ambitious.

% ────────────────────────────────────────────────────────────
\subsection{H1\,---\,Street Canyon Recirculation}

When wind blows perpendicular to a row of buildings (a ``street canyon''), fluid mechanics predicts
that a recirculation vortex forms inside the canyon, trapping pollutants rather than flushing them.
Streets parallel to the wind experience ventilation; perpendicular streets do not. This is one of the
most robust phenomena in urban meteorology and should produce a large, detectable signal even with
modest sampling.

\begin{tcolorbox}[callout, title={Testable Hypothesis}]
PM\textsubscript{2.5} concentrations on streets oriented \textbf{perpendicular} to the prevailing
wind will be 30--50\% higher than on parallel streets of similar traffic density, measured during
steady winds $\geq 3$\,m\,s$^{-1}$.

\medskip
\textbf{Method:} Bin routes by street--wind angle ($0$--$30^\circ$, $31$--$60^\circ$, $61$--$90^\circ$).
One-way ANOVA + Tukey HSD post-hoc test. Report Cohen's $d$ as effect size.\\
\textbf{Control variable:} Annual average daily traffic (AADT) from IDOT open data.
\end{tcolorbox}

% ────────────────────────────────────────────────────────────
\subsection{H2\,---\,Gaussian Plume Dispersion from Highway Line Sources}

% EDIT [physics, 2026-07-22]: AUDIT_papers.md Paper-3 issue #3 -- the substantive
% physics correction in this document.
% WAS: the highway was (correctly) called a continuous line source, but the fall-off
% along a PERPENDICULAR transect was then modelled as a Gaussian in the CROSS-WIND
% coordinate y, C(y) = Q/(sqrt(2pi) sigma_y u) exp(-y^2/2 sigma_y^2), with the fitted
% width interpreted as sigma_y and compared to Pasquill-Gifford curves.
% WHY THAT IS WRONG: for an (effectively) infinite crosswind line source the
% concentration does not fall off in the crosswind direction at all -- by construction
% every crosswind position sees an identical, translation-invariant source. What a
% perpendicular transect actually measures is the decay with DOWNWIND distance x, and
% that decay is governed by VERTICAL dispersion sigma_z(x): the plume dilutes upward,
% not sideways. Fitting the crosswind form to a perpendicular transect returns a number
% that is neither sigma_y nor sigma_z and cannot be compared to any PG curve.
% CORRECTION APPLIED: the standard ground-level infinite-line-source result, in which a
% ground-level receptor at downwind distance x sees
%   C(x) = 2 Q_L / ( sqrt(2/pi) ... ) -> C(x) = sqrt(2/pi) * Q_L / (sigma_z(x) u),
% i.e. concentration falls off purely through the growth of sigma_z(x). Fit
% sigma_z(x) = a x^b and compare the exponent b to the PG class-D sigma_z curve.
% Basis: standard line-source Gaussian plume theory (Turner, Workbook of Atmospheric
% Dispersion Estimates; Pasquill-Gifford sigma_z parameterisation).
% Confidence: high on the diagnosis (sigma_y is unambiguously the wrong coordinate);
% MEDIUM on this specific replacement -- see the FLAG below.
% FLAG FOR AUTHOR: this adopts the INFINITE-line idealisation, valid when the road is
% long compared with the transect distance (fine for I-90/94 out to a few hundred m,
% questionable near interchanges or a road end). The alternative is to treat the
% highway as a FINITE line and integrate along it explicitly, which keeps a genuine
% crosswind dependence and would justify a two-parameter (sigma_y, sigma_z) fit.
% Choose deliberately; the infinite-line form is the standard default and is assumed here.
A highway acts as a continuous line source of PM\textsubscript{2.5} and NO$_x$. Classical
atmospheric dispersion theory (Gaussian plume model) predicts that, for a line source
long compared with the sampling distance, concentration is uniform \emph{along} the road
and falls off with \emph{downwind} distance $x$ through the growth of the vertical
dispersion parameter $\sigma_z(x)$, governed by
Pasquill--Gifford (PG) atmospheric stability classes. (There is no cross-wind fall-off
to fit: an infinite line source is translation-invariant along its own axis, so the
perpendicular-transect decay is a $\sigma_z$ measurement, not a $\sigma_y$ one.)
This is the most directly quantitative
hypothesis---fitting the plume model yields physically interpretable parameters
($\sigma_z$, line emission rate $Q_L$) that can be compared to published values.

\begin{tcolorbox}[callout, title={Testable Hypothesis}]
PM\textsubscript{2.5} concentrations on repeated transects across I-90/94 and I-290 will
decay with downwind distance in a manner consistent with a ground-level line source whose
$\sigma_z(x)$ matches Pasquill--Gifford class D (neutral)
stability at the measured wind speeds.

\medskip
\textbf{Method:} 20+ perpendicular transects per highway. Nonlinear least-squares fit of
the ground-level infinite-line-source relation
\[
  C(x) = \sqrt{\frac{2}{\pi}}\;\frac{Q_L}{\sigma_z(x)\,u},
  \qquad \sigma_z(x) = a\,x^{b}
\]
where $Q_L$ is the emission rate \emph{per unit length of road} and $u$ the
transport wind speed. Extract $\sigma_z(x)$; compare the exponent $b$ and prefactor $a$
to PG class-D $\sigma_z$ power-law curves. Estimate $Q_L$; compare to USEPA NEI emission
factors (per-vehicle factor $\times$ measured traffic count per unit length).
Report the wind angle relative to the road and restrict fits to near-perpendicular
transport, since the line-source form assumes flow across the road rather than along it.
\end{tcolorbox}

% ────────────────────────────────────────────────────────────
\subsection{H3\,---\,Stokes-Number-Dependent Particle Separation}

Coarse particles (PM\textsubscript{10} but not PM\textsubscript{2.5}) have a larger Stokes number
$\mathrm{St} = \tau_p / \tau_f$, where the particle relaxation time is
$\tau_p = \rho_p d_p^2 / 18\mu$. High-St particles cannot follow rapid changes in airflow and
preferentially deposit near their source. Fine particles (PM\textsubscript{2.5}) track the flow
and travel farther. This predicts that the PM\textsubscript{10}/PM\textsubscript{2.5} mass ratio
should decrease with distance from a coarse-particle source and increase with wind speed at fixed
distance.

\begin{tcolorbox}[callout, title={Testable Hypothesis}]
The PM\textsubscript{10}/PM\textsubscript{2.5} mass concentration ratio will decrease
monotonically with distance from identified coarse-particle sources (construction sites, unsealed
lots) at a rate consistent with Stokes-number-predicted settling velocities for 2.5--10\,$\mu$m
diameter particles.

\medskip
\textbf{Method:} Characterize 3--5 point sources. Run 10+ transects per source at distances
20, 50, 100, 200\,m. Fit exponential decay: $R(r) = R_0 \exp(-\lambda r)$. Interpret $\lambda$ via
theoretical settling velocity $v_s = \rho_p d_p^2 g / 18\mu$.
\end{tcolorbox}

% ────────────────────────────────────────────────────────────
\subsection{H4\,---\,Lake Breeze Effects on Near-Surface Pollutant Concentrations}

Lake Michigan drives a mesoscale lake breeze circulation on summer afternoons: cool, relatively
clean air advects inland, deepening the atmospheric boundary layer (ABL) and diluting near-surface
pollutants. This effect should be detectable as a spatial gradient in PM\textsubscript{2.5} along
an east--west transect, conditional on confirmed lake breeze events.

\begin{tcolorbox}[callout, title={Testable Hypothesis}]
On confirmed lake breeze days (easterly surface wind, KMDW vs.\ KORD wind speed divergence),
PM\textsubscript{2.5} concentrations within 1\,km of Lake Michigan will be $\geq 20\%$ lower than
on matched non-lake-breeze days with similar traffic and time-of-day conditions.

\medskip
\textbf{Method:} Composite analysis. Lake breeze days identified via NOAA ISD + surface analysis.
Matched-pair design (same day-of-week, time-of-day, season). Mann--Whitney $U$ test;
effect size $r = Z/\sqrt{N}$.
\end{tcolorbox}

% ────────────────────────────────────────────────────────────
\subsection{Ancillary Goal\,---\,Open Dataset Publication}

Regardless of hypothesis outcomes, the spatial dataset itself is a novel contribution. No open,
street-level, multi-parameter air quality dataset at cycling/running resolution ($\sim$10\,m)
exists for Chicago. The dataset will be published to Zenodo under a CC-BY\,4.0 license with a DOI,
and a Data Descriptor submitted to \textit{Earth System Science Data}.

\newpage

% ============================================================
\section{Hardware Build List}
% ============================================================

Total estimated cost: \textbf{\$165--\$185 USD}. All components are commercially available; no
specialized lab equipment is required. The Sensirion SEN54 is the non-negotiable core sensor---it
combines five measurements into a single factory-calibrated unit with peer-reviewed validation
against federal reference instruments (Zimmerman et al., 2018; Tryner et al., 2020).

\medskip
\begin{center}
\begin{tabular}{L{3.8cm} L{5.0cm} C{1.4cm} L{3.2cm}}
\toprule
\rowcolor{navy!90}
\color{white}\textbf{Component} &
\color{white}\textbf{Notes} &
\color{white}\textbf{Cost} &
\color{white}\textbf{Source} \\
\midrule
Raspberry Pi Zero 2W &
Quad-core 1\,GHz, 512\,MB RAM. Low-power Linux SBC. &
\$15 & Adafruit, PiShop.us \\
\rowcolor{lightgray}
Sensirion SEN54 &
PM\textsubscript{1/2.5/4/10}, VOC, T, RH on one I\textsuperscript{2}C board. Factory-calibrated laser counter. &
\$45 & Mouser, DigiKey \\
u-blox NEO-M9N GPS &
10\,Hz, 2.5\,m CEP. NMEA over UART. Strong urban canyon performance. &
\$30 & SparkFun \#GPS-17285 \\
\rowcolor{lightgray}
MicroSD card, 64\,GB U3 A2 &
Class-10 write speed avoids I/O bottlenecks during 1\,Hz logging. &
\$10 & Amazon (Samsung Endurance) \\
USB-C battery, 10{,}000\,mAh &
Powers Pi + sensors for 8--10\,hrs continuously. &
\$25 & Amazon (Anker PowerCore) \\
\rowcolor{lightgray}
USB-C cable + splitter &
Power routing from pack to Pi. &
\$8 & Amazon \\
JST-SH / Qwiic cable (4-pin) &
I\textsuperscript{2}C connection SEN54 $\to$ Pi GPIO. &
\$5 & SparkFun, Adafruit \\
\rowcolor{lightgray}
CP2102 USB--UART adapter &
UART connection for GPS; doubles as debug serial port. &
\$8 & Amazon \\
ABS enclosure 120$\times$80$\times$55\,mm &
Protects from sweat, rain, condensation. Hammond 1554A or similar. &
\$12 & Amazon \\
\rowcolor{lightgray}
Silica gel packets &
Humidity control inside sealed enclosure. &
\$3 & Amazon \\
40-pin GPIO header &
Required for Pi Zero 2W (ships unsoldered). &
\$3 & Adafruit \\
\rowcolor{lightgray}
Jumper wires, heat shrink, tape &
Wiring and strain relief. &
\$8 & Amazon / local shop \\
3D-printed vest/handlebar mount &
Attach enclosure to running vest or bike bar. STL files on Thingiverse. &
\$5 & Self-print \\
\midrule
\multicolumn{2}{l}{\textbf{Total}} & \textbf{\$177} & \\
\bottomrule
\end{tabular}
\end{center}

\begin{tcolorbox}[callout, title={Why the SEN54 specifically?}]
Sensirion's SEN5x series uses a calibrated laser-diffraction optical particle counter with
on-board signal processing. Independent laboratory evaluations show PM\textsubscript{2.5}
$R^2 > 0.95$ against FEM reference instruments under most atmospheric conditions. The VOC index
is a qualitative relative metric (scale 1--500), not absolute ppb---useful for detecting emission
events but not for regulatory comparison.
\end{tcolorbox}

\newpage

% ============================================================
\section{Step-by-Step Build Instructions}
% ============================================================

\subsection{Hardware Assembly}

\subsubsection{Step 1\,---\,Prepare the Raspberry Pi Zero 2W}
\begin{enumerate}
  \item Solder or press-fit the 40-pin GPIO header if not pre-soldered.
  \item Insert a formatted 64\,GB microSD card (will be overwritten in \S\ref{sec:software}).
  \item Do not power on until software is fully installed.
\end{enumerate}

\subsubsection{Step 2\,---\,Wire the SEN54 via I\textsuperscript{2}C}
The SEN54 communicates over I\textsuperscript{2}C at 3.3\,V logic. Use 4 wires:
\begin{enumerate}
  \item \textbf{VDD} $\to$ Pi GPIO Pin 1 (3.3\,V). SEN54 operates at 3.3\,V---confirmed safe.
  \item \textbf{GND} $\to$ Pi GPIO Pin 6 (GND).
  \item \textbf{SDA} $\to$ Pi GPIO Pin 3 (I\textsuperscript{2}C SDA, BCM GPIO2).
  \item \textbf{SCL} $\to$ Pi GPIO Pin 5 (I\textsuperscript{2}C SCL, BCM GPIO3).
  \item Leave the \texttt{SEL} pin floating or tie to GND (selects I\textsuperscript{2}C mode). Default address: \texttt{0x69}.
  \item Use a Qwiic/JST-SH cable for a secure, rattle-free connection. Protect both ends with heat shrink.
\end{enumerate}

\subsubsection{Step 3\,---\,Wire the GPS Module via UART}
\begin{enumerate}
  \item TX (GPS) $\to$ RX (Pi GPIO Pin 10, BCM GPIO15).
  \item RX (GPS) $\to$ TX (Pi GPIO Pin 8, BCM GPIO14).
  \item VCC (GPS) $\to$ 5\,V or 3.3\,V depending on your module---check the datasheet.
  \item GND (GPS) $\to$ Pi GPIO Pin 14.
  \item \textit{Note:} Disable the Pi's serial console via \texttt{raspi-config} before using UART for GPS.
\end{enumerate}

\subsubsection{Step 4\,---\,Enclosure and Weatherproofing}
\begin{enumerate}
  \item Drill two holes in the ABS lid: one for the SEN54 inlet/outlet ports (must be open to air flow), one for the USB-C power cable.
  \item Use M12 weatherproof cable glands at all entry points.
  \item Place silica gel packets inside to control interior humidity.
  \item Mount Pi, battery, and sensors with double-sided foam tape or M2.5 standoffs.
  \item Test fit: all cables should have slack; no stress on connectors when lid is closed.
  \item Attach to running vest or handlebar with zip ties or outdoor-rated Velcro. Orient so the SEN54 inlet faces into oncoming airflow.
\end{enumerate}

% ────────────────────────────────────────────────────────────
\subsection{Software Setup}
\label{sec:software}

\subsubsection{Step 1\,---\,Install Raspberry Pi OS Lite}
\begin{enumerate}
  \item Download Raspberry Pi OS Lite (64-bit, Bookworm) from \url{https://raspberrypi.com}.
  \item Flash to microSD using Raspberry Pi Imager. In Advanced Options: set hostname, enable SSH, set home WiFi credentials.
  \item Boot, SSH in: \texttt{ssh pi@<hostname>.local}
\end{enumerate}

\subsubsection{Step 2\,---\,Enable Interfaces}
\begin{enumerate}
  \item Run \texttt{sudo raspi-config}
  \item Interface Options $\to$ I\textsuperscript{2}C: \textbf{Enable}.
  \item Interface Options $\to$ Serial Port: disable the login shell; enable the hardware serial port.
  \item Reboot. Verify with \texttt{i2cdetect -y 1}---you should see \texttt{0x69} (SEN54).
\end{enumerate}

\subsubsection{Step 3\,---\,Install Python Dependencies}
\begin{enumerate}
  \item \texttt{sudo apt update \&\& sudo apt install -y python3-pip python3-serial gpsd gpsd-clients}
  \item \texttt{pip3 install sensirion-i2c-driver sensirion-i2c-sen5x pynmea2}
  \item \texttt{pip3 install pandas pyserial}
\end{enumerate}

\subsubsection{Step 4\,---\,Configure \texttt{gpsd}}
\begin{enumerate}
  \item Edit \texttt{/etc/default/gpsd}: set \texttt{DEVICES="/dev/ttyS0"} and \texttt{GPSD\_OPTIONS="-n"}.
  \item \texttt{sudo systemctl enable gpsd \&\& sudo systemctl start gpsd}
  \item Test: \texttt{cgps -s}---should show lat/lon within 60\,s in open sky.
\end{enumerate}

\subsubsection{Step 5\,---\,Deploy the Logging Script}

Create \texttt{/home/pi/airlog/logger.py}. The core loop reads sensor and GPS at 1\,Hz and
writes one row per second to a date-stamped CSV:

\begin{tcolorbox}[callout, title={Logger Script---Core Logic (Python pseudocode)}]
\small\ttfamily
sensor = Sen5xI2cDevice(I2cConnection(i2c\_transceiver))\\
sensor.start\_measurement()\\
gps\_ser = serial.Serial('/dev/ttyS0', 9600, timeout=1)\\
\\
with open(f'/data/airlog\_\{date.today()\}.csv', 'w') as f:\\
\quad writer.writerow(['timestamp','lat','lon','pm1','pm25',\\
\quad\quad\quad\quad\quad\quad\quad\quad\quad\quad 'pm4','pm10','voc\_index','temp\_c',\\
\quad\quad\quad\quad\quad\quad\quad\quad\quad\quad 'rh\_pct','fix\_quality','n\_sats'])\\
\quad while True:\\
\quad\quad parse \$GNGGA NMEA sentence $\to$ lat, lon, fix, sats\\
\quad\quad read sensor $\to$ pm1, pm25, pm4, pm10, rh, temp, voc\\
\quad\quad write row with datetime.utcnow().isoformat()\\
\quad\quad time.sleep(1)   \# 1\,Hz sampling rate
\end{tcolorbox}

\subsubsection{Step 6\,---\,Auto-Start on Boot}
\begin{enumerate}
  \item Create \texttt{/etc/systemd/system/airlog.service} with \texttt{ExecStart} pointing to \texttt{logger.py}.
  \item \texttt{sudo systemctl enable airlog \&\& sudo systemctl start airlog}
  \item Reboot; verify CSVs appear in \texttt{/data/} within 30\,s.
\end{enumerate}

\subsubsection{Step 7\,---\,Pre-Run Calibration}
\begin{itemize}
  \item Co-locate the SEN54 beside a PurpleAir or IQAir AirVisual for 24\,hrs. Record any systematic PM\textsubscript{2.5} offset. If offset $> 15\%$, document a correction factor.
  \item Verify GPS cold-start time $< 90$\,s in open sky.
  \item Log a 10-minute stationary test. Verify timestamps are UTC, lat/lon stable, all columns populated.
\end{itemize}

\newpage

% ============================================================
\section{Data Collection Protocol}
% ============================================================

\subsection{Route Design Strategy}

Routes are designed to isolate specific independent variables, not chosen at random. Build four
route families:

\subsubsection{Route Family 1\,---\,Street Orientation Transects (for H1)}
\begin{itemize}
  \item Design 3--5 matched pairs: one block running N--S, one E--W, within 200\,m of each other with comparable building height and traffic density.
  \item Target Chicago's grid: Wacker, Michigan Ave, State St vs.\ adjacent cross-streets.
  \item Run each pair on the same day within 30\,min so meteorological conditions are identical.
  \item \textbf{Go/no-go condition:} steady wind $\geq 3$\,m\,s$^{-1}$. Check ASOS KMDW before leaving.
\end{itemize}

\subsubsection{Route Family 2\,---\,Highway Plume Transects (for H2)}
\begin{itemize}
  \item Target two highways: I-90/94 (Dan Ryan Expressway) and I-290 (Eisenhower Expressway).
  \item Design perpendicular transects at 25, 50, 100, 200, and 400\,m downwind from the highway centerline.
  \item A transect is only valid when wind is blowing from the highway toward you (confirmed from KMDW obs.).
  \item Run consistently at 7--9\,AM (peak traffic) for comparability across dates.
\end{itemize}

\subsubsection{Route Family 3\,---\,Lakefront Gradient (for H4)}
\begin{itemize}
  \item Design an E--W transect from the lakefront (e.g., Montrose Beach, 41.96$^\circ$\,N) inland $\sim 5$\,km.
  \item Collect data on lake breeze days and non-lake-breeze days; label using NOAA ISD wind records.
  \item Minimum 10 lake breeze runs and 10 matched control runs for adequate comparison.
\end{itemize}

\subsubsection{Route Family 4\,---\,Point Source Transects (for H3)}
\begin{itemize}
  \item Identify 3--5 active construction sites or unsealed lots. Confirm still active before each run.
  \item Run at 20, 50, 100, and 200\,m from the source boundary on concentric arcs.
  \item Target steady moderate winds ($\geq 4$\,m\,s$^{-1}$, from source toward runner).
\end{itemize}

% ────────────────────────────────────────────────────────────
\subsection{Flagging Wind and Weather Conditions}

For every run, record the following alongside the sensor data:

\medskip
\begin{center}
\begin{tabular}{L{3.2cm} L{3.8cm} L{6.0cm}}
\toprule
\rowcolor{navy!90}
\color{white}\textbf{Variable} &
\color{white}\textbf{Source} &
\color{white}\textbf{How to Record} \\
\midrule
Wind speed (10\,m) &
NOAA ASOS KMDW or KORD &
Iowa State Mesonet API; log obs.\ at start and end of run \\
\rowcolor{lightgray}
Wind direction &
Same ASOS station &
Bin into 8 cardinal directions. Flag runs with variable wind ($>45^\circ$ std.\ dev.) \\
Stability class &
Wind speed + cloud cover + solar elevation &
Pasquill--Gifford table (Seinfeld \& Pandis, Table 4.2) \\
\rowcolor{lightgray}
Precipitation &
NWS observed &
Exclude runs during or within 2\,hrs of rain (wet deposition removes PM) \\
Mixing layer height &
NOAA HRRR model (hourly) &
AWS Open Data. Low MLH ($\leq 500$\,m) = pollution trapping event \\
\rowcolor{lightgray}
Lake breeze flag &
KMDW vs.\ KORD wind divergence &
Binary yes/no per run; confidence notes in field log \\
Traffic proxy &
Time of day + day of week &
AM peak (7--9), midday (10--15), PM peak (16--19), off-peak \\
\bottomrule
\end{tabular}
\end{center}

% ────────────────────────────────────────────────────────────
\subsection{Statistical Power and Required Sample Size}

Estimates assume $\alpha = 0.05$, power $= 0.80$, effect sizes from published urban air quality
literature.

\medskip
\begin{center}
\begin{tabular}{L{3.0cm} L{3.8cm} L{2.8cm} C{2.0cm}}
\toprule
\rowcolor{navy!90}
\color{white}\textbf{Hypothesis} &
\color{white}\textbf{Expected Effect} &
\color{white}\textbf{Runs Needed} &
\color{white}\textbf{Running Total} \\
\midrule
H1: Street canyon &
Cohen's $d \approx 0.8$, 30--50\% PM\textsubscript{2.5} difference &
$\sim$60 route segments &
60 \\
\rowcolor{lightgray}
H2: Highway plume &
Gaussian fit $R^2 > 0.7$ &
40 transects (2 highways) &
100 \\
H3: Particle separation &
Moderate effect, high variability &
40 transects (4 sources) &
140 \\
\rowcolor{lightgray}
H4: Lake breeze &
20--30\% PM reduction &
20 runs &
160 \\
Buffer / exploratory &
Opportunistic repeats &
$+40$ runs &
\textbf{200} \\
\bottomrule
\end{tabular}
\end{center}

\medskip
\textbf{Bottom line:} $\sim$200 total logged runs over 6--12 months. At 3--4 runs per week
(integrated into your normal exercise habit), this is achievable in 50--67 weeks. Prioritize H1
and H2 first---they are weather-gated and cannot be rescheduled freely.

\newpage

% ============================================================
\section{Data Pipeline}
% ============================================================

The pipeline transforms raw, messy SD-card CSVs into a clean, analysis-ready dataset. Every step
should be scripted and fully reproducible. Use Python with \texttt{pandas}, \texttt{geopandas},
and \texttt{pyarrow}.

\subsubsection{Step 1\,---\,Ingest and Concatenate}
\begin{itemize}
  \item Read all per-run CSV files. Tag each row with \texttt{run\_id}, \texttt{route\_family}, and \texttt{run\_date}.
  \item Assert the column schema on read---fail loudly on any missing column rather than silently dropping data.
  \item Save concatenated raw data to \texttt{raw\_combined.parquet} (preserves dtypes, much faster than CSV for repeated reads).
\end{itemize}

\subsubsection{Step 2\,---\,GPS Quality Filter}
\begin{itemize}
  \item Drop rows where \texttt{fix\_quality\,==\,0} (no GPS fix) or \texttt{n\_sats\,<\,4}.
  \item Flag rows where implied speed between consecutive positions exceeds physically plausible values ($>35$\,km\,h$^{-1}$ running, $>60$\,km\,h$^{-1}$ cycling).
  \item Apply a Hampel identifier (5-point median filter, threshold $= 3\sigma$) to \texttt{lat} and \texttt{lon} separately to remove spike outliers.
  \item Report fraction of data dropped per run; flag any run losing $>10\%$ of rows.
\end{itemize}

\subsubsection{Step 3\,---\,Sensor Quality Control}
\begin{itemize}
  \item PM\textsubscript{2.5}, PM\textsubscript{10}: physical range 0--1000\,$\mu$g\,m$^{-3}$. Flag out-of-range values as instrument errors.
  \item Temperature: flag values outside $-20$ to $50^\circ$C.
  \item RH $>85\%$: flag as high-humidity condition. Optical particle counters overcount PM\textsubscript{2.5} at high RH due to hygroscopic particle growth; exclude these rows from plume fitting or analyze separately.
  \item VOC index: qualitative scale 1--500. Use only for event detection, not absolute comparison.
  \item Sensor warm-up artifact: SEN54 requires $\sim$30\,s to stabilize. Drop the first 60\,s of every run.
\end{itemize}

\subsubsection{Step 4\,---\,Time Synchronization}
\begin{itemize}
  \item Verify Pi clock was NTP-synced before each run (add NTP check to logger startup script).
  \item GPS NMEA sentences carry UTC time---use this as the master clock, not the Pi system clock.
  \item Interpolate GPS position to 1\,Hz if GPS dropouts create sub-second gaps.
\end{itemize}

\subsubsection{Step 5\,---\,Meteorological Data Merge}
\begin{itemize}
  \item Download hourly ASOS observations for KMDW and KORD from the Iowa State Mesonet API (\url{https://mesonet.agron.iastate.edu}).
  \item For each run, assign the ASOS observation nearest in time ($\pm 30$\,min) as the meteorological context.
  \item Compute PG stability class: wind speed + cloud cover + solar elevation angle (Seinfeld \& Pandis, Eq.~4.7).
\end{itemize}

\subsubsection{Step 6\,---\,Spatial Binning}
\begin{itemize}
  \item Project GPS coordinates to Illinois State Plane (EPSG:3435) for accurate metric distance calculations.
  \item Snap each measurement to a $10 \times 10$\,m grid cell using a \texttt{geopandas} spatial join.
  \item Per cell compute: mean PM\textsubscript{2.5}, std PM\textsubscript{2.5}, $n_\text{obs}$, $n_\text{runs}$. Flag cells with $n_\text{obs} < 5$ as low-confidence.
\end{itemize}

\subsubsection{Step 7\,---\,Final Dataset Schema}

Two output tables:
\begin{itemize}
  \item \texttt{observations.parquet}: one row per 1-second sensor reading, all quality flags applied.
  \item \texttt{grid\_summary.parquet}: one row per 10\,m grid cell, aggregated statistics across all runs.
\end{itemize}

\medskip
\begin{center}
\small
\begin{tabular}{L{3.8cm} L{3.2cm} L{6.0cm}}
\toprule
\rowcolor{navy!90}
\color{white}\textbf{Column} &
\color{white}\textbf{Type} &
\color{white}\textbf{Description} \\
\midrule
\texttt{run\_id} & \texttt{str} & Unique ID per run (\texttt{YYYYMMDD\_routeID}) \\
\rowcolor{lightgray}
\texttt{timestamp\_utc} & \texttt{datetime64[ns,UTC]} & GPS-anchored UTC timestamp \\
\texttt{lat} / \texttt{lon} & \texttt{float64} & WGS84 decimal degrees, Hampel-filtered \\
\rowcolor{lightgray}
\texttt{easting} / \texttt{northing} & \texttt{float64} & Illinois State Plane (m) \\
\texttt{pm1}, \texttt{pm25}, \texttt{pm4}, \texttt{pm10} & \texttt{float32} & Particle mass concentration ($\mu$g\,m$^{-3}$) \\
\rowcolor{lightgray}
\texttt{voc\_index} & \texttt{int16} & Sensirion VOC index (1--500, dimensionless) \\
\texttt{temp\_c}, \texttt{rh\_pct} & \texttt{float32} & Ambient temperature ($^\circ$C) and relative humidity (\%) \\
\rowcolor{lightgray}
\texttt{wind\_speed\_ms}, \texttt{wind\_dir\_deg} & \texttt{float32} & From nearest ASOS obs., interpolated to run time \\
\texttt{stability\_class} & \texttt{str} & Pasquill--Gifford A--F \\
\rowcolor{lightgray}
\texttt{lake\_breeze\_flag} & \texttt{bool} & \texttt{True} if run tagged as lake breeze event \\
\texttt{route\_family} & \texttt{categorical} & \texttt{canyon / plume / lakefront / point\_source / exploratory} \\
\rowcolor{lightgray}
\texttt{qc\_flag} & \texttt{int8} & Bitmask: bit0=GPS drop, bit1=high RH, bit2=warm-up, bit3=range \\
\bottomrule
\end{tabular}
\end{center}

\newpage

% ============================================================
\section{Analysis Roadmap}
\label{sec:analysis}
% ============================================================

Work through these levels in order. Each one builds on the previous and serves as a sanity check
that the data are behaving physically before you commit to model fitting.

% ────────────────────────────────────────────────────────────
\subsection{Level 1\,---\,Exploratory Spatial Analysis}
\begin{itemize}
  \item \textbf{Spatial heatmap:} Plot all PM\textsubscript{2.5} measurements on an interactive Chicago basemap using \texttt{folium} or \texttt{kepler.gl}. This is your first look at whether the data make physical sense---hot spots should correspond to highways, industrial zones, and street canyons, not random locations.
  \item \textbf{Time-of-day profiles:} Boxplot PM\textsubscript{2.5} by hour of day. Expect morning and evening peaks corresponding to commuter traffic; clean midday if winds are strong.
  \item \textbf{Seasonal decomposition:} Compare summer vs.\ winter runs. Winter temperature inversions trap boundary-layer pollutants and should produce dramatically higher concentrations than summer runs on the same routes.
  \item \textbf{Sensor self-consistency:} Scatter PM\textsubscript{2.5} vs.\ PM\textsubscript{10}. Expect $R^2 > 0.7$. Points far off the regression line (anomalously high PM\textsubscript{10}/PM\textsubscript{2.5}) signal coarse-particle events worth investigating.
  \item \textbf{Wind rose of polluted runs:} 90th-percentile PM\textsubscript{2.5} binned by wind sector. A quick qualitative check on dominant source directions before any model fitting.
\end{itemize}

% ────────────────────────────────────────────────────────────
\subsection{Level 2\,---\,Statistical Testing of Core Hypotheses}
\begin{itemize}
  \item \textbf{H1 (Street canyon):} One-way ANOVA across street--wind angle bins. Tukey HSD post-hoc test. Report Cohen's $d$ as effect size. Visualize as mean $\pm$ 95\% CI bar chart.
  \item \textbf{H4 (Lake breeze):} Mann--Whitney $U$ test (non-parametric; PM distributions are right-skewed). Effect size $r = Z/\sqrt{N}$.
  \item \textbf{For both:} Bootstrap confidence intervals ($n = 10{,}000$ resamples) on group means; this requires no distributional assumptions and is straightforward to implement with \texttt{scipy.stats} or \texttt{numpy}.
\end{itemize}

% ────────────────────────────────────────────────────────────
\subsection{Level 3\,---\,Gaussian Plume Fitting (H2)}
\begin{itemize}
% EDIT [physics, 2026-07-22]: analysis pipeline brought into line with the corrected H2
% above. The old recipe binned by "cross-wind distance from the highway centerline" and
% fitted a Gaussian in y -- but distance from the centerline along a perpendicular
% transect IS the downwind coordinate x, not the crosswind coordinate y. The two were
% conflated. Refitting the decay through sigma_z(x) as the line-source model requires.
% Note this also removes a latent bug: fitting exp(-y^2/2 sigma^2) to a monotonically
% decaying transect (no peak inside the sampled range beyond the roadside maximum)
% is poorly conditioned and curve_fit would often return an arbitrary width.
% Confidence: high. See the flag at the H2 hypothesis for the infinite- vs finite-line
% modelling choice.
  \item For each highway transect, bin PM\textsubscript{2.5} by \emph{downwind} distance $x$ from the highway edge (10\,m bins). Record the wind angle relative to the road and keep only near-perpendicular transport.
  \item Fit the ground-level line-source decay
  \[
    C(x) = \frac{C_0}{\sigma_z(x)}, \qquad \sigma_z(x) = a\,x^{b}
    \quad\Longrightarrow\quad C(x) = C_0\,a^{-1} x^{-b}
  \]
  using \texttt{scipy.optimize.curve\_fit} (fit in log--log space, where the model is linear and the fit is well conditioned). Extract $a$, $b$ and their uncertainties.
  \item Compare the exponent $b$ to PG stability-class $\sigma_z$ predictions (class D neutral: $b \approx 0.7$--$0.8$ over the 10--300\,m range; verify against the specific PG parameterisation used).
  \item Estimate the per-unit-length emission rate $Q_L$ from the fitted prefactor via $C(x) = \sqrt{2/\pi}\,Q_L/(\sigma_z(x)u)$ with measured $u$. Convert to a per-vehicle factor using the traffic count and compare to USEPA National Emissions Inventory (NEI) highway emission factors.
  \item \textbf{Background subtraction:} fit $C(x) = C_\mathrm{bg} + C_0 a^{-1}x^{-b}$, or subtract a control-site background first --- a power-law decay fitted without a background term will bias $b$ low.
\end{itemize}

% ────────────────────────────────────────────────────────────
\subsection{Level 4\,---\,Stokes Number Analysis (H3)}
\begin{itemize}
  \item Compute particle relaxation time: $\tau_p = \rho_p d_p^2 / 18\mu$ for $d_p = 2.5$ and $10\,\mu$m. Use $\rho_p = 1500$\,kg\,m$^{-3}$ (typical urban aerosol), $\mu = 1.81\times10^{-5}$\,Pa\,s.
  \item Compute $\mathrm{St} = \tau_p U / L$ where $U$ is measured wind speed and $L$ is a characteristic building scale ($\sim$20\,m).
  \item Plot PM\textsubscript{10}/PM\textsubscript{2.5} ratio vs.\ distance from each point source. Fit exponential decay $R(r) = R_0 e^{-\lambda r}$.
  \item Interpret $\lambda$ via theoretical settling velocity $v_s = \rho_p d_p^2 g / 18\mu$. \textit{This level is the most original physics in the project---in-situ Stokes analysis at street scale is genuinely novel in the literature.}
\end{itemize}

% ────────────────────────────────────────────────────────────
\subsection{Level 5\,---\,Data-Driven Spatial Modeling}
\begin{itemize}
  \item Train a random forest regressor to predict PM\textsubscript{2.5} at any (lat, lon) given: wind speed, wind direction, time of day, day of week, proximity to roads/highways (OpenStreetMap network), land use (NLCD), and season.
  \item Evaluate with RMSE and spatial cross-validation (leave-one-neighborhood-out), to prevent inflated performance from spatial autocorrelation.
  \item Extract feature importances to quantify which variables dominate PM\textsubscript{2.5} predictions across Chicago.
  \item Compare to a kriging interpolation (geostatistical baseline).
\end{itemize}

% ────────────────────────────────────────────────────────────
\subsection{Level 6\,---\,Community Impact Map}
\begin{itemize}
  \item Join \texttt{grid\_summary} with the Chicago Community Areas shapefile (City of Chicago Open Data Portal).
  \item Compute mean and 90th-percentile PM\textsubscript{2.5} per community area.
  \item Overlay with demographic data (ACS 5-year estimates: \% below poverty line, \% non-white residents, \% under 5 and over 65 years).
  \item Report Pearson correlation between socioeconomic indicators and mean PM\textsubscript{2.5}. This is the public-facing story and the most immediately communicable result.
\end{itemize}

\newpage

% ============================================================
\section{Milestone Timeline (6--12 Months)}
% ============================================================

Designed to be completed alongside a full PhD workload---roughly 5--10 hours per week of
project time (including the runs themselves).

\medskip
\begin{center}
\begin{tabular}{L{2.0cm} L{2.4cm} L{9.0cm}}
\toprule
\rowcolor{navy!90}
\color{white}\textbf{Period} &
\color{white}\textbf{Phase} &
\color{white}\textbf{Key Deliverables and Actions} \\
\midrule
Month 1 &
Hardware &
Order all components. Assemble Pi + SEN54 + GPS. Verify I\textsuperscript{2}C and UART. Log first 30-min stationary test file. Confirm GPS cold-start time. \\
\rowcolor{lightgray}
Months 1--2 &
Software &
Deploy \texttt{logger.py} with systemd auto-start. Build pre-run health check script. 24-hr co-location calibration vs.\ PurpleAir or IQAir. Document any correction factors. \\
Months 2--3 &
Pilot Runs &
20 exploratory runs across all route families to shake out hardware and logging issues. Build v1 data pipeline. Fix systematic errors. Refine enclosure weatherproofing. \\
\rowcolor{lightgray}
Months 3--6 &
Systematic Collection &
\textbf{Priority phase.} Focus on weather-gated H1 and H2 runs. Target 3--4 runs/week. Reach $\sim$100 runs by month 6. Build met.\ merge pipeline in parallel. \\
Months 6--7 &
Pipeline \& QC &
Run full data pipeline on all collected data. Generate QC report (flag drop rates, RH exceedances, GPS quality). Produce first spatial heatmap. Identify gaps; plan fill runs. \\
\rowcolor{lightgray}
Months 7--8 &
Exploratory Analysis &
Complete Level 1 and Level 2 analyses. Run all 4 hypothesis tests. Share informal update with advisor or lab group. Refine collection strategy based on power analysis of real data. \\
Months 8--10 &
Advanced Analysis &
Gaussian plume fitting (Level 3). Stokes analysis (Level 4). ML spatial model (Level 5). Freeze dataset for publication. \\
\rowcolor{lightgray}
Months 10--11 &
Write-Up &
Write dataset descriptor paper (\textit{ESSD} target). Write main results paper or technical report. Produce blog post and GitHub repository. \\
Months 11--12 &
Dissemination &
Submit dataset to Zenodo. Post to personal site. Submit abstract to AGU or AAAR. Consider journal submission (see \S\ref{sec:present}). \\
\bottomrule
\end{tabular}
\end{center}

\begin{tcolorbox}[callout, title={PhD Compatibility Notes}]
\begin{itemize}[topsep=0pt]
  \item Data collection = 45--90\,min per run, integrated into exercise you would do anyway.
  \item Analysis is modular: each Level in \S\ref{sec:analysis} is a self-contained Jupyter notebook.
  \item Hardware is passive once deployed---no babysitting required during runs.
  \item Build pipeline and analysis during slower PhD periods (between semesters, post-quals).
  \item A 3-month delay in any phase is recoverable without compromising scientific quality.
\end{itemize}
\end{tcolorbox}

\newpage

% ============================================================
\section{How to Present the Final Result}
\label{sec:present}
% ============================================================

\subsection{GitHub Repository (Public Portfolio Piece)}
\begin{itemize}
  \item \textbf{Repository structure:} \texttt{/hardware} (wiring diagrams, enclosure CAD), \texttt{/firmware} (logger.py, systemd config), \texttt{/pipeline} (data cleaning notebooks), \texttt{/analysis} (one notebook per hypothesis), \texttt{/data} (Zenodo links), \texttt{/docs} (rendered GitHub Pages site).
  \item Write a thorough \texttt{README.md} with photos of the hardware setup, example output maps, and one-sentence plain-language summaries of each finding.
  \item Use DVC (Data Version Control) or Git LFS to track large data files without bloating the repo.
  \item \textbf{Target:} full reproducibility---a collaborator should be able to \texttt{git clone}, download the Zenodo dataset, run \texttt{make analysis}, and reproduce every figure.
\end{itemize}

\subsection{Dataset Publication}
\begin{itemize}
  \item Publish the clean dataset to Zenodo (\url{https://zenodo.org}) under CC-BY\,4.0. Free, issues a DOI, indexed by Google Scholar and OpenAIRE.
  \item Consider a Data Descriptor submission to \textit{Earth System Science Data} (ESSD)---a peer-reviewed EGU journal for novel environmental datasets (IF $\approx 12$). An ESSD paper is a standalone, citable publication even if it precedes the main analysis paper.
  \item Minimum metadata: sensor make/model/firmware version, calibration protocol and correction factors, GPS accuracy characterization, route family definitions, and meteorological data sources.
\end{itemize}

\subsection{Blog Post / Public Communication}
\begin{itemize}
  \item Write one 2{,}000-word accessible post for a technically curious general audience. Suggested platforms: Towards Data Science, a personal Hugo/Jekyll site, or Substack.
  \item Suggested structure: (1) Chicago's monitoring gap, (2) the hardware you built and why, (3) one striking map or chart, (4) what the physics tells us, (5) what you'll do next.
  \item Cross-post on Twitter/X, LinkedIn, and Reddit (\texttt{r/dataisbeautiful}, \texttt{r/chicago}, \texttt{r/STEM}). This is how citizen science projects get noticed by journalists and environmental agencies.
\end{itemize}

\subsection{Conference Presentations}
\begin{itemize}
  \item \textbf{AAAR Annual Conference} (American Association for Aerosol Research): ideal venue. Abstract deadline typically April for October conference. Target session: Urban Aerosols / Low-Cost Sensors.
  \item \textbf{AGU Fall Meeting}: large atmospheric science community. Target session: Atmospheric Composition / Air Quality.
  \item \textbf{AMS Annual Meeting}: good fit for the lake breeze and Gaussian dispersion components.
  \item \textbf{Northwestern Graduate Research Symposium}: low-stakes entry point to practice your talk. Present there first.
\end{itemize}

\subsection{Potential Journal Venues}

\medskip
\begin{center}
\begin{tabular}{L{4.2cm} C{1.8cm} L{7.2cm}}
\toprule
\rowcolor{navy!90}
\color{white}\textbf{Journal} &
\color{white}\textbf{Fit} &
\color{white}\textbf{Notes} \\
\midrule
\textit{Atmospheric Measurement Techniques} (EGU) &
Excellent &
Best fit if sensor validation + novel dataset is the primary contribution. Open access. \\
\rowcolor{lightgray}
\textit{Atmospheric Environment} &
Excellent &
Best fit if main findings are the urban aerosol physics results (H1--H4). Broad readership. \\
\textit{npj Climate and Atmospheric Science} (Nature) &
Strong &
High-impact. Requires novel, broadly significant findings. Good if equity/lake breeze results are striking. \\
\rowcolor{lightgray}
\textit{Environmental Science \& Technology} (ACS) &
Good &
Well-regarded applied urban air quality journal; strong policy-community readership. \\
\textit{Earth System Science Data} (EGU) &
Excellent (dataset) &
Peer-reviewed data descriptor. Citable standalone publication; does not require full analysis. \\
\rowcolor{lightgray}
\textit{PLOS ONE} &
Fallback &
Accepts well-executed work regardless of novelty threshold. Good if hypothesis results are mixed. \\
\bottomrule
\end{tabular}
\end{center}

\begin{tcolorbox}[callout, title={Final Advice}]
\textbf{Start writing before the analysis is done.} Write the Introduction and Methods sections while you are still collecting data---the Methods section is fully knowable right now.

\medskip
\textbf{The dataset is the guaranteed deliverable.} Even if every hypothesis yields a null result, a publicly released, GPS-tagged, street-level air quality dataset for Chicago is a genuine scientific contribution.

\medskip
\textbf{This project demonstrates, in one package:} hardware design, experimental protocol, data engineering, statistical analysis, and physical modeling. For a physics PhD, that range of demonstrated skills is rare and impressive to future employers in academia, national labs, and industry alike.
\end{tcolorbox}

% ============================================================
\end{document}
% ============================================================
